
\documentclass{beamer}

\usetheme{Madrid}
\usecolortheme{seahorse}
\usepackage{fontspec}
\usepackage{polyglossia}
\setmainlanguage{english} 
\setotherlanguage{russian}
\setotherlanguage{chinese}



\setmainfont{Times New Roman}  
\setsansfont{Arial}  

\setbeamercolor{frametitle}{bg=blue!70!black, fg=white}
\setbeamerfont{frametitle}{series=\bfseries, size=\large}


\usepackage{graphicx} 

\begin{document}

\begin{frame}{Раздел 6: Развернутые системы LLM}
\begin{itemize}
    \item Разработка LLM сталкивается с проблемами, но выявляет стратегии их преодоления.
    \item Два направления:
    \begin{enumerate}
        \item Предобучение: повышение устойчивости к сбоям с помощью диагностики и автоматического восстановления.
        \item Оценка: быстрая реакция на задачи через декомпозицию.
    \end{enumerate}
\end{itemize}
\end{frame}


\begin{frame}{6.1 Устойчивое к сбоям предобучение}
\textbf{Мотивация:} Предобучение LLM сопровождается сбоями из-за большого числа GPU и длительности обучения. Это снижает эффективность использования ресурсов и требует ручного вмешательства инженеров.

\textbf{Система:} Автоматическое выявление причин сбоев и восстановление включает три модуля:
\begin{enumerate}
    \item \textbf{Чекпоинтинг:} частое сохранение модели с использованием асинхронного подхода для снижения накладных расходов.
    \item \textbf{Диагностика:} комбинация правил и LLM для точного выявления причин сбоев.
    \item \textbf{Восстановление:} обнаружение проблемных узлов и автоматический рестарт с последнего чекпоинта.
\end{enumerate}
\end{frame}


\begin{frame}{Диагностика сбоев и восстановление модели}
\begin{center}
\includegraphics[width=0.9\textwidth,height=0.7\textheight,keepaspectratio]{failure_workflow.png}
\end{center}
\begin{itemize}
    \item Рисунок демонстрирует поток обработки логов, фильтрацию и восстановление.
    \item Система значительно снижает ручное вмешательство (около 90\%).
\end{itemize}
\end{frame}


\begin{frame}{6.2 Разделённое планирование оценивания}
\textbf{Мотивация:} Оценка качества модели по одному метрику не даёт полной картины. Необходимо:
\begin{itemize}
    \item Регулярная оценка на множестве задач для отслеживания прогресса.
    \item Быстрая обратная связь для корректировки обучения.
\end{itemize}

\textbf{Система:} Координатор задач реализует:
\begin{enumerate}
    \item Разделение загрузки модели с удалённого хранилища и процесса оценивания.
    \item Разделение вычисления метрик на отдельные CPU задачи.
    \item Приоритетное планирование на основе ожидаемой длительности задач.
\end{enumerate}
\end{frame}


\begin{frame}{Нагрузочное тестирование и планирование оценивания}
\begin{center}
\includegraphics[width=0.9\textwidth,height=0.7\textheight,keepaspectratio]{valuation_scheduling.png}
\end{center}
\begin{itemize}
    \item Левая часть: При 8 задачах загрузка падает из-за NIC — узкое место.
    \item Правая часть: Модель в CPU-кэш, оценка через PCIe — без повторной загрузки, GPU не ждёт.
    \item Декомпозиция повышает эффективность и сокращает время простоя GPU.
\end{itemize}
\end{frame}